\newpage
\begin{center}
  {\large\bfseries ПРИЛОЖЕНИЕ А}\\[8pt]
  \textbf{(обязательное)}\\[8pt]
  {\large\bfseries ЛИСТИНГ ПРОГРАММЫ}
\end{center}
\refstepcounter{chapter}
\addcontentsline{toc}{chapter}{ПРИЛОЖЕНИЕ А. Листинг программы}

\lstset{
  basicstyle=\ttfamily\footnotesize,
  breaklines=true,
  breakatwhitespace=false,
  columns=fullflexible,
  keepspaces=true,
  showstringspaces=false,
  frame=single,
  tabsize=4,
  language=Python,
}

\bigskip
\noindent{\large\bfseries \_\_init\_\_.py}
\nopagebreak
\begin{lstlisting}
"""PyTorch PINN implementation for the rising bubble two-phase flow case."""

from .config import TrainingConfig, preset_config
from .model import TwoPhasePINN

__all__ = ["TrainingConfig", "TwoPhasePINN", "preset_config"]
\end{lstlisting}

\bigskip
\noindent{\large\bfseries config.py}
\nopagebreak
\begin{lstlisting}
from __future__ import annotations

from dataclasses import asdict, dataclass, field
from pathlib import Path


@dataclass
class PointConfig:
    alpha: tuple[int, int] = (120, 100)
    pde: tuple[int, int, int] = (100, 300, 400)
    north: tuple[int, int] = (12, 12)
    south: tuple[int, int] = (12, 12)
    east: tuple[int, int] = (12, 12)
    west: tuple[int, int] = (12, 12)


@dataclass
class PhysicsConfig:
    mu: tuple[float, float] = (1.0, 10.0)
    rho: tuple[float, float] = (100.0, 1000.0)
    sigma: float = 24.5
    g: float = -0.98
    u_ref: float = 1.0
    l_ref: float = 0.25
    p_ref_rho: float = 1000.0


@dataclass
class TrainingConfig:
    data_path: Path = Path("cfd_data/rising_bubble.h5")
    output_dir: Path = Path("diploma/checkpoints")
    hidden_layers: tuple[int, ...] = (128, 128, 128, 128)
    activation: str = "tanh"
    points: PointConfig = field(default_factory=PointConfig)
    physics: PhysicsConfig = field(default_factory=PhysicsConfig)
    loss_weights_pde: tuple[float, float, float, float] = (1.0, 10.0, 10.0, 1.0)
    parameterized: bool = True
    epochs: tuple[int, ...] = (1000, 1000, 1000)
    learning_rates: tuple[float, ...] = (1e-4, 5e-5, 1e-5)
    batch_size: int | None = 4096
    num_batches: int | None = None
    checkpoint_interval: int = 100
    seed: int = 1234
    device: str = "auto"
    dtype: str = "float32"
    num_workers: int = 0

    def to_dict(self) -> dict:
        data = asdict(self)
        data["data_path"] = str(self.data_path)
        data["output_dir"] = str(self.output_dir)
        return data


def preset_config(name: str) -> TrainingConfig:
    if name == "smoke":
        return TrainingConfig(
            hidden_layers=(32, 32),
            points=PointConfig(
                alpha=(20, 20),
                pde=(20, 40, 40),
                north=(6, 4),
                south=(6, 4),
                east=(6, 4),
                west=(6, 4),
            ),
            epochs=(2,),
            learning_rates=(1e-3,),
            batch_size=512,
            checkpoint_interval=1,
        )
    if name == "default":
        return TrainingConfig()
    if name == "paper":
        return TrainingConfig(
            hidden_layers=(350,) * 8,
            points=PointConfig(
                alpha=(500, 400),
                pde=(400, 2000, 3000),
                north=(20, 20),
                south=(20, 20),
                east=(20, 20),
                west=(20, 20),
            ),
            epochs=(5000,) * 5,
            learning_rates=(1e-4, 5e-5, 1e-5, 5e-6, 1e-6),
            batch_size=None,
            num_batches=20,
            checkpoint_interval=100,
        )
    if name == "paper_light":
        return TrainingConfig(
            hidden_layers=(32,64,128,256,128,64,32),
            points=PointConfig(
                alpha=(500, 400),
                pde=(400, 2000, 3000),
                north=(20, 20),
                south=(20, 20),
                east=(20, 20),
                west=(20, 20),
            ),
            epochs=(5000,) * 5,
            learning_rates=(1e-4, 5e-5, 1e-5, 5e-6, 1e-6),
            batch_size=None,
            num_batches=20,
            checkpoint_interval=100,
        )
    raise ValueError(f"Unknown preset '{name}'. Use smoke, default, paper_light, or paper.")
\end{lstlisting}

\bigskip
\noindent{\large\bfseries model.py}
\nopagebreak
\begin{lstlisting}
from __future__ import annotations

from dataclasses import dataclass

import torch
from torch import nn
import torch.nn.functional as F

from .config import PhysicsConfig


def _activation(name: str) -> nn.Module:
    if name == "tanh":
        return nn.Tanh()
    if name == "silu":
        return nn.SiLU()
    if name == "gelu":
        return nn.GELU()
    if name == "sin":
        return Sine()
    raise ValueError(f"Unsupported activation '{name}'.")


class Sine(nn.Module):
    def forward(self, x: torch.Tensor) -> torch.Tensor:
        return torch.sin(x)


class MLP(nn.Module):
    def __init__(self, input_dim: int, hidden_layers: tuple[int, ...], activation: str = "tanh"):
        super().__init__()
        layers: list[nn.Module] = []
        previous = input_dim
        for width in hidden_layers:
            linear = nn.Linear(previous, width)
            nn.init.xavier_normal_(linear.weight)
            nn.init.zeros_(linear.bias)
            layers.extend([linear, _activation(activation)])
            previous = width

        self.trunk = nn.Sequential(*layers)
        self.u_head = nn.Linear(previous, 1)
        self.v_head = nn.Linear(previous, 1)
        self.p_head = nn.Linear(previous, 1)
        self.alpha_head = nn.Linear(previous, 1)
        for head in (self.u_head, self.v_head, self.p_head, self.alpha_head):
            nn.init.xavier_normal_(head.weight)
            nn.init.zeros_(head.bias)

    def forward(self, xyt: torch.Tensor) -> tuple[torch.Tensor, torch.Tensor, torch.Tensor, torch.Tensor]:
        z = self.trunk(xyt)
        u = self.u_head(z)
        v = self.v_head(z)
        p = torch.exp(torch.clamp(self.p_head(z), min=-20.0, max=20.0))
        alpha = torch.sigmoid(self.alpha_head(z))
        return u, v, p, alpha


@dataclass
class LossBreakdown:
    total: torch.Tensor
    alpha: torch.Tensor
    boundary: torch.Tensor
    mass: torch.Tensor
    momentum_x: torch.Tensor
    momentum_y: torch.Tensor
    alpha_pde: torch.Tensor

    def detached(self) -> dict[str, float]:
        return {
            "total": float(self.total.detach().cpu()),
            "alpha": float(self.alpha.detach().cpu()),
            "boundary": float(self.boundary.detach().cpu()),
            "mass": float(self.mass.detach().cpu()),
            "momentum_x": float(self.momentum_x.detach().cpu()),
            "momentum_y": float(self.momentum_y.detach().cpu()),
            "alpha_pde": float(self.alpha_pde.detach().cpu()),
        }


class TwoPhasePINN(nn.Module):
    def __init__(
        self,
        hidden_layers: tuple[int, ...],
        physics: PhysicsConfig,
        activation: str = "tanh",
        loss_weights_pde: tuple[float, float, float, float] = (1.0, 10.0, 10.0, 1.0),
        input_dim: int = 4,
    ):
        super().__init__()
        self.input_dim = input_dim
        self.net = MLP(input_dim, hidden_layers, activation)
        self.physics = physics
        self.register_buffer("loss_weights_pde", torch.tensor(loss_weights_pde, dtype=torch.float32))

    def forward(self, xyt: torch.Tensor) -> tuple[torch.Tensor, torch.Tensor, torch.Tensor, torch.Tensor]:
        return self.net(xyt)

    @staticmethod
    def _grad(y: torch.Tensor, x: torch.Tensor) -> torch.Tensor:
        return torch.autograd.grad(
            y,
            x,
            grad_outputs=torch.ones_like(y),
            retain_graph=True,
            create_graph=True,
            only_inputs=True,
        )[0]

    def pde_residuals(self, xyt: torch.Tensor) -> tuple[torch.Tensor, torch.Tensor, torch.Tensor, torch.Tensor]:
        xyt = xyt.detach().clone().requires_grad_(True)
        u, v, p, alpha = self.forward(xyt)

        grad_u = self._grad(u, xyt)
        grad_v = self._grad(v, xyt)
        grad_p = self._grad(p, xyt)
        grad_alpha = self._grad(alpha, xyt)

        u_x, u_y, u_t = grad_u[:, 0:1], grad_u[:, 1:2], grad_u[:, 2:3]
        v_x, v_y, v_t = grad_v[:, 0:1], grad_v[:, 1:2], grad_v[:, 2:3]
        p_x, p_y = grad_p[:, 0:1], grad_p[:, 1:2]
        a_x, a_y, a_t = grad_alpha[:, 0:1], grad_alpha[:, 1:2], grad_alpha[:, 2:3]

        u_xx = self._grad(u_x, xyt)[:, 0:1]
        u_yy = self._grad(u_y, xyt)[:, 1:2]
        v_xx = self._grad(v_x, xyt)[:, 0:1]
        v_yy = self._grad(v_y, xyt)[:, 1:2]
        a_xx = self._grad(a_x, xyt)[:, 0:1]
        a_yy = self._grad(a_y, xyt)[:, 1:2]
        a_xy = self._grad(a_x, xyt)[:, 1:2]

        phys = self.physics
        mu1, mu2 = phys.mu
        rho1, rho2 = phys.rho
        mu = mu2 + (mu1 - mu2) * alpha
        mu_x = (mu1 - mu2) * a_x
        mu_y = (mu1 - mu2) * a_y
        rho = rho2 + (rho1 - rho2) * alpha

        eps = torch.finfo(xyt.dtype).eps
        abs_interface_grad = torch.sqrt(a_x.square() + a_y.square() + eps)
        curvature = -(
            (a_xx + a_yy) / abs_interface_grad
            - (a_x.square() * a_xx + a_y.square() * a_yy + 2.0 * a_x * a_y * a_xy)
            / abs_interface_grad.pow(3)
        )

        rho_ref = rho2
        one_re = mu / (rho_ref * phys.u_ref * phys.l_ref)
        one_re_x = mu_x / (rho_ref * phys.u_ref * phys.l_ref)
        one_re_y = mu_y / (rho_ref * phys.u_ref * phys.l_ref)
        one_we = phys.sigma / (rho_ref * phys.u_ref**2 * phys.l_ref)
        one_fr = phys.g * phys.l_ref / phys.u_ref**2

        mass = u_x + v_y
        alpha_transport = a_t + u * a_x + v * a_y
        momentum_x = (
            (u_t + u * u_x + v * u_y) * rho / rho_ref
            + p_x
            - one_we * curvature * a_x
            - one_re * (u_xx + u_yy)
            - 2.0 * one_re_x * u_x
            - one_re_y * (u_y + v_x)
        )
        momentum_y = (
            (v_t + u * v_x + v * v_y) * rho / rho_ref
            + p_y
            - one_we * curvature * a_y
            - one_re * (v_xx + v_yy)
            - rho / rho_ref * one_fr
            - 2.0 * one_re_y * v_y
            - one_re_x * (u_y + v_x)
        )
        return mass, momentum_x, momentum_y, alpha_transport

    def loss(self, batches: dict[str, torch.Tensor]) -> LossBreakdown:
        input_dim = self.input_dim
        alpha_batch = batches["alpha"]
        _, _, _, pred_alpha = self.forward(alpha_batch[:, :input_dim])
        loss_alpha = F.mse_loss(pred_alpha, alpha_batch[:, input_dim : input_dim + 1])

        north = batches["north"]
        _, _, p_n, _ = self.forward(north[:, :input_dim])
        loss_north_p = F.mse_loss(p_n, north[:, input_dim : input_dim + 1])

        east_west = batches["east_west"]
        u_e, v_e, p_e, _ = self.forward(east_west[:, :input_dim])
        u_w, v_w, p_w, _ = self.forward(east_west[:, input_dim : 2 * input_dim])
        loss_ew = F.mse_loss(u_e, u_w) + F.mse_loss(v_e, v_w) + F.mse_loss(p_e, p_w)

        nsew = batches["nsew"]
        u_b, v_b, _, _ = self.forward(nsew[:, :input_dim])
        loss_velocity_bc = F.mse_loss(u_b, nsew[:, input_dim : input_dim + 1]) + F.mse_loss(
            v_b, nsew[:, input_dim + 1 : input_dim + 2]
        )
        loss_boundary = loss_north_p + loss_ew + loss_velocity_bc

        pde_batch = batches["pde"]
        target = pde_batch[:, input_dim : input_dim + 1]
        mass, momentum_x, momentum_y, alpha_transport = self.pde_residuals(pde_batch[:, :input_dim])
        loss_mass = F.mse_loss(mass, target)
        loss_momentum_x = F.mse_loss(momentum_x, target)
        loss_momentum_y = F.mse_loss(momentum_y, target)
        loss_alpha_pde = F.mse_loss(alpha_transport, target)

        weights = self.loss_weights_pde.to(loss_mass.device, dtype=loss_mass.dtype)
        loss_pde = (
            weights[0] * loss_mass
            + weights[1] * loss_momentum_x
            + weights[2] * loss_momentum_y
            + weights[3] * loss_alpha_pde
        )
        total = loss_alpha + loss_boundary + loss_pde
        return LossBreakdown(
            total=total,
            alpha=loss_alpha,
            boundary=loss_boundary,
            mass=loss_mass,
            momentum_x=loss_momentum_x,
            momentum_y=loss_momentum_y,
            alpha_pde=loss_alpha_pde,
        )

    @torch.no_grad()
    def predict(self, xyt: torch.Tensor, batch_size: int = 1_000_000) -> tuple[torch.Tensor, ...]:
        outputs = []
        for chunk in xyt.split(batch_size):
            outputs.append(self.forward(chunk))
        return tuple(torch.cat([item[i] for item in outputs], dim=0) for i in range(4))
\end{lstlisting}

\bigskip
\noindent{\large\bfseries data.py}
\nopagebreak
\begin{lstlisting}
from __future__ import annotations

from dataclasses import dataclass
from pathlib import Path
import math
import re
from collections.abc import Sequence

import h5py
import numpy as np
import torch

from .config import PointConfig


@dataclass
class TrainingData:
    alpha: torch.Tensor
    pde: torch.Tensor
    north: torch.Tensor
    east_west: torch.Tensor
    nsew: torch.Tensor

    def to(self, device: torch.device, dtype: torch.dtype) -> "TrainingData":
        return TrainingData(
            alpha=self.alpha.to(device=device, dtype=dtype),
            pde=self.pde.to(device=device, dtype=dtype),
            north=self.north.to(device=device, dtype=dtype),
            east_west=self.east_west.to(device=device, dtype=dtype),
            nsew=self.nsew.to(device=device, dtype=dtype),
        )

    @property
    def sizes(self) -> dict[str, int]:
        return {
            "alpha": len(self.alpha),
            "pde": len(self.pde),
            "north": len(self.north),
            "east_west": len(self.east_west),
            "nsew": len(self.nsew),
        }


def _choice(rng: np.random.Generator, values: np.ndarray, size: int, replace: bool = False) -> np.ndarray:
    if size <= len(values):
        return rng.choice(values, size, replace=replace)
    return rng.choice(values, size, replace=True)


def _points_for_interface(
    rng: np.random.Generator,
    n_points: int,
    t_value: float,
    interface: list[tuple[float, float]],
    normal: list[tuple[float, float]],
    tangent: list[tuple[float, float]],
    cell_size: float,
    refine_start: float,
    refine_end: float,
) -> np.ndarray:
    if not interface:
        raise ValueError("No interface cells were detected in the selected CFD snapshot.")

    per_cell = math.ceil(n_points / (2 * len(interface)))
    points_inward: list[np.ndarray] = []
    points_outward: list[np.ndarray] = []

    for x_p, n_p, t_p in zip(interface, normal, tangent):
        x_p = np.asarray(x_p, dtype=np.float64)
        n_p = np.asarray(n_p, dtype=np.float64)
        t_p = np.asarray(t_p, dtype=np.float64)
        inward_n = rng.uniform(refine_start, refine_end, per_cell)
        outward_n = rng.uniform(refine_start, refine_end, per_cell)
        inward_t1 = rng.uniform(refine_start, refine_end, per_cell)
        inward_t2 = rng.uniform(refine_start, refine_end, per_cell)
        outward_t1 = rng.uniform(refine_start, refine_end, per_cell)
        outward_t2 = rng.uniform(refine_start, refine_end, per_cell)

        for in_n, in_t1, in_t2, out_n, out_t1, out_t2 in zip(
            inward_n, inward_t1, inward_t2, outward_n, outward_t1, outward_t2
        ):
            points_inward.append(x_p + in_n * n_p)
            points_inward.append(x_p + cell_size / 3.0 * t_p + in_t1 * n_p)
            points_inward.append(x_p + 2.0 * cell_size / 3.0 * t_p + in_t2 * n_p)
            points_outward.append(x_p - out_n * n_p)
            points_outward.append(x_p + cell_size / 3.0 * t_p - out_t1 * n_p)
            points_outward.append(x_p + 2.0 * cell_size / 3.0 * t_p - out_t2 * n_p)

    inward = np.hstack(
        [
            np.asarray(points_inward),
            t_value * np.ones((len(points_inward), 1)),
            np.ones((len(points_inward), 1)),
        ]
    )
    outward = np.hstack(
        [
            np.asarray(points_outward),
            t_value * np.ones((len(points_outward), 1)),
            np.zeros((len(points_outward), 1)),
        ]
    )
    data = np.vstack([inward, outward])
    return data[_choice(rng, np.arange(len(data)), n_points), :]


def _points_for_domain(
    rng: np.random.Generator,
    n_points: int,
    t_value: float,
    x: np.ndarray,
    y: np.ndarray,
    levelset: np.ndarray,
    cell_size: float,
    max_levelset: float,
) -> np.ndarray:
    per_x = math.ceil(n_points / len(x))
    data = np.empty((0, 4), dtype=np.float64)
    levelset_threshold = max_levelset * cell_size

    for index_x, x_value in enumerate(x):
        t_domain = t_value * np.ones((per_x, 1))
        x_domain = x_value * np.ones((per_x, 1))
        y_chunks: list[np.ndarray] = []
        a_chunks: list[np.ndarray] = []
        remaining = per_x

        while remaining:
            indices_y = _choice(rng, np.arange(len(y)), remaining)
            y_temp = y[indices_y]
            levelset_temp = levelset[index_x, indices_y]
            a_temp = (levelset_temp > 0).astype(np.float64)
            keep = np.where(np.abs(levelset_temp) >= levelset_threshold)[0]
            if len(keep) == 0:
                continue
            y_chunks.append(y_temp[keep])
            a_chunks.append(a_temp[keep])
            remaining -= len(keep)

        y_domain = np.hstack(y_chunks).reshape(per_x, 1)
        a_domain = np.hstack(a_chunks).reshape(per_x, 1)
        data = np.vstack([data, np.hstack([x_domain, y_domain, t_domain, a_domain])])

    return data[_choice(rng, np.arange(len(data)), n_points), :]


def _compute_normals(
    x: np.ndarray, y: np.ndarray, levelset: np.ndarray, cell_size: float
) -> tuple[list, list, list]:
    interfaces: list[list[tuple[float, float]]] = []
    normals: list[list[tuple[float, float]]] = []
    tangents: list[list[tuple[float, float]]] = []

    for levelset_t in levelset:
        ix, iy = np.where(np.abs(levelset_t) <= 0.75 * cell_size)
        valid = (ix > 0) & (ix < len(x) - 1) & (iy > 0) & (iy < len(y) - 1)
        ix = ix[valid]
        iy = iy[valid]

        interfaces.append(list(zip(x[ix], y[iy])))
        ls_x = (levelset_t[ix + 1, iy] - levelset_t[ix - 1, iy]) / (2.0 * cell_size)
        ls_y = (levelset_t[ix, iy + 1] - levelset_t[ix, iy - 1]) / (2.0 * cell_size)
        grad_abs = np.sqrt(ls_x**2 + ls_y**2)
        grad_abs[grad_abs == 0.0] = np.finfo(float).eps
        normal_x = ls_x / grad_abs
        normal_y = ls_y / grad_abs
        normals.append(list(zip(normal_x, normal_y)))
        tangents.append(list(zip(-normal_y, normal_x)))

    return interfaces, normals, tangents


def _points_alpha(
    rng: np.random.Generator,
    cfg: tuple[int, int],
    times: np.ndarray,
    x: np.ndarray,
    y: np.ndarray,
    cell_size: float,
    interfaces: list,
    normals: list,
    tangents: list,
    levelset: np.ndarray,
) -> np.ndarray:
    chunks = []
    for vals in zip(times, interfaces, normals, tangents, levelset):
        t_value, interface, normal, tangent, levelset_t = vals
        chunks.append(_points_for_interface(rng, cfg[0], t_value, interface, normal, tangent, cell_size, 0.004, 0.008))
        chunks.append(_points_for_domain(rng, cfg[1], t_value, x, y, levelset_t, cell_size, 4.0))
    return np.vstack(chunks)


def _points_pde(
    rng: np.random.Generator,
    cfg: tuple[int, int, int],
    times: np.ndarray,
    x: np.ndarray,
    y: np.ndarray,
    cell_size: float,
    interfaces: list,
    normals: list,
    tangents: list,
    levelset: np.ndarray,
) -> np.ndarray:
    chunks = []
    for vals in zip(times, interfaces, normals, tangents, levelset):
        t_value, interface, normal, tangent, levelset_t = vals
        chunks.append(_points_for_interface(rng, cfg[0], t_value, interface, normal, tangent, cell_size, 0.0, 0.001))
        chunks.append(_points_for_interface(rng, cfg[1], t_value, interface, normal, tangent, cell_size, 0.001, 0.1))
        chunks.append(_points_for_domain(rng, cfg[2], t_value, x, y, levelset_t, cell_size, 4.0))
    data = np.vstack(chunks)
    data[:, 3] = 0.0
    return data


def _north(rng: np.random.Generator, cfg: tuple[int, int], x_bounds, y_bounds, t_bounds) -> np.ndarray:
    low_t = t_bounds[0] + np.finfo(float).eps
    times = np.hstack([0.0, rng.uniform(low_t, t_bounds[1], max(cfg[1] - 2, 0)), t_bounds[1]])
    data = np.empty((0, 6), dtype=np.float64)
    for t_value in times:
        xs = rng.uniform(x_bounds[0], x_bounds[1], (cfg[0], 1))
        ys = y_bounds[1] * np.ones((cfg[0], 1))
        data = np.vstack([data, np.hstack([xs, ys, t_value * np.ones((cfg[0], 1)), np.zeros((cfg[0], 2)), np.ones((cfg[0], 1))])])
    return data


def _wall(cfg: tuple[int, int], x_bounds, y_bounds, t_bounds, side: str) -> np.ndarray:
    times = np.linspace(t_bounds[0], t_bounds[1], cfg[1])
    data = np.empty((0, 5), dtype=np.float64)
    for t_value in times:
        if side == "south":
            xs = np.linspace(x_bounds[0], x_bounds[1], cfg[0]).reshape(-1, 1)
            ys = y_bounds[0] * np.ones((cfg[0], 1))
        elif side == "east":
            xs = x_bounds[1] * np.ones((cfg[0], 1))
            ys = np.linspace(y_bounds[0], y_bounds[1], cfg[0]).reshape(-1, 1)
        elif side == "west":
            xs = x_bounds[0] * np.ones((cfg[0], 1))
            ys = np.linspace(y_bounds[0], y_bounds[1], cfg[0]).reshape(-1, 1)
        else:
            raise ValueError(side)
        data = np.vstack([data, np.hstack([xs, ys, t_value * np.ones((cfg[0], 1)), np.zeros((cfg[0], 2))])])
    return data


def selected_time_indices(n_times: int) -> np.ndarray:
    indices = np.arange(n_times)
    return np.sort(np.concatenate([indices[0:30:15], indices[30:100:5], indices[100::5], indices[1:3]], axis=0))


def resolve_data_paths(data_path: Path | Sequence[Path]) -> list[Path]:
    if isinstance(data_path, (str, Path)):
        path = Path(data_path)
        if path.is_dir():
            paths = sorted(path.glob("*.h5"))
            if not paths:
                raise FileNotFoundError(f"No .h5 files found in directory: {path}")
            return paths
        return [path]
    return [Path(path) for path in data_path]


def radius_from_h5_path(data_path: Path, default: float = 0.25) -> float:
    with h5py.File(data_path, "r") as data:
        for key in ("radius", "R"):
            if key in data.attrs:
                return float(data.attrs[key])

    match = re.search(r"(?:^|[_-])R(\d+(?:\.\d+)?)", data_path.stem, flags=re.IGNORECASE)
    if match:
        value = match.group(1)
        radius = float(value)
        if radius >= 1.0 and "." not in value:
            radius /= 100.0
        return radius

    return default


def _append_radius_column(data: np.ndarray, radius: float, target_start: int = 3) -> np.ndarray:
    radius_col = radius * np.ones((len(data), 1), dtype=data.dtype)
    return np.hstack([data[:, :target_start], radius_col, data[:, target_start:]])


def _concat_training_data(chunks: list[TrainingData]) -> TrainingData:
    return TrainingData(
        alpha=torch.cat([chunk.alpha for chunk in chunks], dim=0),
        pde=torch.cat([chunk.pde for chunk in chunks], dim=0),
        north=torch.cat([chunk.north for chunk in chunks], dim=0),
        east_west=torch.cat([chunk.east_west for chunk in chunks], dim=0),
        nsew=torch.cat([chunk.nsew for chunk in chunks], dim=0),
    )


def make_training_data(
    data_path: Path | Sequence[Path],
    cfg: PointConfig,
    l_ref: float,
    seed: int = 1234,
    parameterized: bool = True,
) -> TrainingData:
    data_paths = resolve_data_paths(data_path)
    if not parameterized and len(data_paths) != 1:
        raise ValueError("Ordinary PINN training requires exactly one HDF5 data file.")
    chunks = [
        _make_training_data_single(path, cfg, l_ref, seed + index, parameterized)
        for index, path in enumerate(data_paths)
    ]
    return _concat_training_data(chunks)


def _make_training_data_single(
    data_path: Path,
    cfg: PointConfig,
    l_ref: float,
    seed: int,
    parameterized: bool,
) -> TrainingData:
    data_path = Path(data_path)
    if not data_path.exists():
        raise FileNotFoundError(
            f"CFD file not found: {data_path}. Download rising_bubble.h5 and pass --data-path if needed."
        )

    rng = np.random.default_rng(seed)
    radius = radius_from_h5_path(data_path) / l_ref
    with h5py.File(data_path, "r") as data:
        x = np.asarray(data["X"])
        y = np.asarray(data["Y"])
        times = np.asarray(data["time"])
        levelset = -np.asarray(data["levelset"])

    indices = selected_time_indices(len(times))
    times = times[indices]
    levelset = levelset[indices]
    x_bounds = (x[0], x[-1])
    y_bounds = (y[0], y[-1])
    t_bounds = (times[0], times[-1])
    cell_size = float(np.diff(x)[0])

    interfaces, normals, tangents = _compute_normals(x, y, levelset, cell_size)
    alpha = _points_alpha(rng, cfg.alpha, times, x, y, cell_size, interfaces, normals, tangents, levelset)
    pde = _points_pde(rng, cfg.pde, times, x, y, cell_size, interfaces, normals, tangents, levelset)
    north = _north(rng, cfg.north, x_bounds, y_bounds, t_bounds)
    south = _wall(cfg.south, x_bounds, y_bounds, t_bounds, "south")
    east = _wall(cfg.east, x_bounds, y_bounds, t_bounds, "east")
    west = _wall(cfg.west, x_bounds, y_bounds, t_bounds, "west")

    nsew_coords = np.vstack([north[:, 0:3], south[:, 0:3], east[:, 0:3], west[:, 0:3]])
    nsew_for_pde = np.hstack([nsew_coords, np.zeros((len(nsew_coords), 1))])

    for arr in (alpha, pde, north, south, east, west, nsew_coords, nsew_for_pde):
        arr[:, :3] /= l_ref

    nsew_raw = np.vstack([north[:, 0:5], south, east[: cfg.east[0]], west[: cfg.west[0]]])

    if parameterized:
        alpha = _append_radius_column(alpha, radius)
        pde = _append_radius_column(pde, radius)
        nsew_for_pde = _append_radius_column(nsew_for_pde, radius)
        east_west = np.hstack(
            [
                east[:, 0:3],
                radius * np.ones((len(east), 1), dtype=east.dtype),
                west[:, 0:3],
                radius * np.ones((len(west), 1), dtype=west.dtype),
            ]
        )
        nsew = _append_radius_column(nsew_raw, radius)
        north_out = _append_radius_column(north[:, [0, 1, 2, 5]], radius)
    else:
        east_west = np.hstack([east[:, 0:3], west[:, 0:3]])
        nsew = nsew_raw
        north_out = north[:, [0, 1, 2, 5]]

    pde = np.vstack([pde, nsew_for_pde])

    return TrainingData(
        alpha=torch.from_numpy(alpha.astype(np.float32)),
        pde=torch.from_numpy(pde.astype(np.float32)),
        north=torch.from_numpy(north_out.astype(np.float32)),
        east_west=torch.from_numpy(east_west.astype(np.float32)),
        nsew=torch.from_numpy(nsew.astype(np.float32)),
    )


def load_cfd(data_path: Path, start: int = 0, end: int = 151, temporal_step: int = 10, spatial_step: int = 2):
    with h5py.File(data_path, "r") as data:
        return {
            "x": np.asarray(data["X"])[::spatial_step],
            "y": np.asarray(data["Y"])[::spatial_step],
            "time": np.asarray(data["time"])[start:end:temporal_step],
            "levelset": np.asarray(data["levelset"])[start:end:temporal_step, ::spatial_step, ::spatial_step],
            "pressure": np.asarray(data["pressure"])[start:end:temporal_step, ::spatial_step, ::spatial_step],
            "velocity_x": np.asarray(data["velocityX"])[start:end:temporal_step, ::spatial_step, ::spatial_step],
            "velocity_y": np.asarray(data["velocityY"])[start:end:temporal_step, ::spatial_step, ::spatial_step],
        }
\end{lstlisting}

\bigskip
\noindent{\large\bfseries train.py}
\nopagebreak
\begin{lstlisting}
from __future__ import annotations

import argparse
import json
from pathlib import Path
import random
import time

import numpy as np
import torch
from tqdm import trange

from .config import TrainingConfig, preset_config
from .data import TrainingData, make_training_data
from .model import TwoPhasePINN


def resolve_device(device: str) -> torch.device:
    if device == "auto":
        return torch.device("cuda" if torch.cuda.is_available() else "cpu")
    requested = torch.device(device)
    if requested.type == "cuda" and not torch.cuda.is_available():
        raise RuntimeError("CUDA was requested, but torch.cuda.is_available() is False.")
    return requested


def resolve_dtype(dtype: str) -> torch.dtype:
    if dtype == "float32":
        return torch.float32
    if dtype == "float64":
        return torch.float64
    raise ValueError("dtype must be float32 or float64")


def set_seed(seed: int) -> None:
    random.seed(seed)
    np.random.seed(seed)
    torch.manual_seed(seed)
    if torch.cuda.is_available():
        torch.cuda.manual_seed_all(seed)


def make_run_dir(output_dir: Path) -> Path:
    timestamp = time.strftime("%Y%m%d_%H%M%S")
    run_dir = output_dir / timestamp
    run_dir.mkdir(parents=True, exist_ok=False)
    return run_dir


def compute_total_batch_size(data: TrainingData, batch_size: int | None, num_batches: int | None) -> int:
    sizes = data.sizes
    total = sum(sizes.values())
    if batch_size is not None:
        return batch_size
    if num_batches is None:
        return total
    if num_batches <= 0:
        raise ValueError("num_batches must be a positive integer.")
    return int(np.ceil(total / num_batches))


def dataset_batch_sizes(data: TrainingData, total_batch_size: int) -> dict[str, int]:
    sizes = data.sizes
    total = sum(sizes.values())
    if total_batch_size < 0 or total_batch_size >= total:
        return sizes
    return {key: max(1, int(np.ceil(total_batch_size * value / total))) for key, value in sizes.items()}


def iter_batches(data: TrainingData, batch_sizes: dict[str, int]):
    sizes = data.sizes
    number_of_batches = max(int(np.ceil(sizes[key] / batch_sizes[key])) for key in sizes)
    permutations = {key: torch.randperm(sizes[key], device=getattr(data, key).device) for key in sizes}
    for batch_no in range(number_of_batches):
        batch = {}
        for key in sizes:
            tensor = getattr(data, key)
            start = batch_no * batch_sizes[key]
            end = min(start + batch_sizes[key], sizes[key])
            if start >= sizes[key]:
                idx = torch.randint(0, sizes[key], (batch_sizes[key],), device=tensor.device)
            else:
                idx = permutations[key][start:end]
            batch[key] = tensor[idx]
        yield batch


def save_checkpoint(
    path: Path,
    model: TwoPhasePINN,
    optimizer: torch.optim.Optimizer,
    epoch: int,
    stage: int,
    loss: float,
    history: list[dict[str, float]],
    cfg: TrainingConfig,
) -> None:
    torch.save(
        {
            "model_state": model.state_dict(),
            "optimizer_state": optimizer.state_dict(),
            "epoch": epoch,
            "stage": stage,
            "loss": loss,
            "history": history,
            "config": cfg.to_dict(),
        },
        path,
    )


class TrainingRun:
    def __init__(self, cfg: TrainingConfig, run_dir: Path | None = None):
        if len(cfg.epochs) != len(cfg.learning_rates):
            raise ValueError("epochs and learning_rates must have the same length.")

        self.cfg = cfg
        set_seed(cfg.seed)
        self.device = resolve_device(cfg.device)
        self.dtype = resolve_dtype(cfg.dtype)

        print(f"Loading and sampling training data from {cfg.data_path}")
        self.data = make_training_data(
            cfg.data_path,
            cfg.points,
            cfg.physics.l_ref,
            cfg.seed,
            parameterized=cfg.parameterized,
        ).to(self.device, self.dtype)
        print("Training point counts:", self.data.sizes)

        self.run_dir = make_run_dir(cfg.output_dir) if run_dir is None else Path(run_dir)
        self.run_dir.mkdir(parents=True, exist_ok=True)
        if run_dir is None:
            (self.run_dir / "config.json").write_text(json.dumps(cfg.to_dict(), indent=2), encoding="utf-8")

        input_dim = 4 if cfg.parameterized else 3
        self.model = TwoPhasePINN(
            cfg.hidden_layers,
            cfg.physics,
            cfg.activation,
            cfg.loss_weights_pde,
            input_dim=input_dim,
        ).to(
            device=self.device,
            dtype=self.dtype,
        )
        self.optimizer = torch.optim.Adam(self.model.parameters(), lr=cfg.learning_rates[0])
        self.total_batch_size = compute_total_batch_size(self.data, cfg.batch_size, cfg.num_batches)
        self.batch_sizes = dataset_batch_sizes(self.data, self.total_batch_size)
        self.number_of_batches = max(
            int(np.ceil(self.data.sizes[key] / self.batch_sizes[key])) for key in self.data.sizes
        )
        print(
            f"Device: {self.device}, dtype: {self.dtype}, "
            f"total batch size: {self.total_batch_size}, batches/epoch: {self.number_of_batches}, "
            f"batch sizes: {self.batch_sizes}"
        )

        self.history: list[dict[str, float]] = []
        self.best_loss = float("inf")
        self.global_epoch = 0
        self.current_stage = 0

    def load_checkpoint(self, checkpoint_path: Path) -> None:
        checkpoint = torch.load(checkpoint_path, map_location=self.device)
        self.model.load_state_dict(checkpoint["model_state"])
        self.optimizer.load_state_dict(checkpoint["optimizer_state"])
        self.history = checkpoint.get("history", [])
        self.global_epoch = int(checkpoint.get("epoch", 0))
        self.current_stage = int(checkpoint.get("stage", 0))
        self.best_loss = min((row["total"] for row in self.history), default=float(checkpoint.get("loss", "inf")))
        print(
            f"Resumed from {checkpoint_path}: "
            f"epoch={self.global_epoch}, stage={self.current_stage}, best_loss={self.best_loss:.3e}"
        )

    def _write_history(self) -> None:
        (self.run_dir / "history.json").write_text(json.dumps(self.history, indent=2), encoding="utf-8")

    def _save_checkpoint(self, name: str, loss: float, stage: int) -> None:
        save_checkpoint(
            self.run_dir / name,
            self.model,
            self.optimizer,
            self.global_epoch,
            stage,
            loss,
            self.history,
            self.cfg,
        )

    def train_stage(
        self,
        epochs: int | None = None,
        lr: float | None = None,
        stage: int | None = None,
        progress: bool = True,
    ) -> list[dict[str, float]]:
        if stage is None:
            stage = self.current_stage + 1
        if epochs is None:
            epochs = self.cfg.epochs[stage - 1]
        if lr is None:
            lr = self.cfg.learning_rates[stage - 1]

        for group in self.optimizer.param_groups:
            group["lr"] = lr

        stage_history: list[dict[str, float]] = []
        epoch_iter = trange(epochs, desc=f"stage {stage} lr={lr:g}") if progress else range(epochs)
        for _ in epoch_iter:
            self.global_epoch += 1
            epoch_losses: dict[str, float] = {}
            batches = 0

            for batch in iter_batches(self.data, self.batch_sizes):
                self.optimizer.zero_grad(set_to_none=True)
                loss = self.model.loss(batch)
                loss.total.backward()
                torch.nn.utils.clip_grad_norm_(self.model.parameters(), max_norm=100.0)
                self.optimizer.step()

                detached = loss.detached()
                for key, value in detached.items():
                    epoch_losses[key] = epoch_losses.get(key, 0.0) + value
                batches += 1

            averaged = {key: value / batches for key, value in epoch_losses.items()}
            averaged["epoch"] = self.global_epoch
            averaged["stage"] = stage
            averaged["lr"] = lr
            self.history.append(averaged)
            stage_history.append(averaged)

            if progress:
                epoch_iter.set_postfix(
                    total=f"{averaged['total']:.3e}",
                    alpha=f"{averaged['alpha']:.2e}",
                    pde_u=f"{averaged['momentum_x']:.2e}",
                )

            if averaged["total"] < self.best_loss:
                self.best_loss = averaged["total"]
                self._save_checkpoint("best.pt", self.best_loss, stage)

            if self.cfg.checkpoint_interval and self.global_epoch % self.cfg.checkpoint_interval == 0:
                self._save_checkpoint(f"epoch_{self.global_epoch:06d}.pt", averaged["total"], stage)
                self._write_history()

        self.current_stage = max(self.current_stage, stage)
        self._save_checkpoint("last.pt", self.history[-1]["total"], stage)
        self._write_history()
        return stage_history

    def train_all(self) -> Path:
        start_stage = self.current_stage + 1
        for stage, (stage_epochs, lr) in enumerate(zip(self.cfg.epochs, self.cfg.learning_rates), start=start_stage):
            self.train_stage(stage_epochs, lr, stage=stage)
        print(f"Finished. Best checkpoint: {self.run_dir / 'best.pt'}")
        return self.run_dir


def train(cfg: TrainingConfig, resume: Path | None = None) -> Path:
    run = TrainingRun(cfg, run_dir=resume.parent if resume is not None else None)
    if resume is not None:
        run.load_checkpoint(resume)
    return run.train_all()


def parse_args() -> argparse.Namespace:
    parser = argparse.ArgumentParser(description="Train a PyTorch PINN for the rising bubble case.")
    parser.add_argument("--preset", choices=["smoke", "default", "paper_light", "paper"], default="default")
    parser.add_argument("--data-path", type=Path)
    parser.add_argument("--output-dir", type=Path)
    parser.add_argument("--device", default=None, help="auto, cpu, cuda, cuda:0, ...")
    parser.add_argument("--dtype", choices=["float32", "float64"], default=None)
    parser.add_argument("--epochs", type=int, default=None, help="Override with a single training stage.")
    parser.add_argument("--batch-size", type=int, default=None)
    parser.add_argument("--num-batches", type=int, default=None, help="Compute total batch size as ceil(total_points / num_batches).")
    parser.add_argument("--hidden-width", type=int, default=None)
    parser.add_argument("--hidden-depth", type=int, default=None)
    parser.add_argument("--lr", type=float, default=None, help="Override with a single learning rate stage.")
    parser.add_argument("--resume", type=Path, default=None, help="Resume from a checkpoint such as checkpoints/.../last.pt.")
    parser.add_argument("--seed", type=int, default=None)
    parser.add_argument(
        "--ordinary-pinn",
        action="store_false",
        dest="parameterized",
        help="Train a non-parametric model with inputs (x, y, t). Requires a single HDF5 data file.",
    )
    parser.add_argument(
        "--parameterized-pinn",
        action="store_true",
        dest="parameterized",
        help="Train a parametric model with inputs (x, y, t, R).",
    )
    parser.set_defaults(parameterized=None)
    return parser.parse_args()


def main() -> None:
    args = parse_args()
    cfg = preset_config(args.preset)
    if args.data_path is not None:
        cfg.data_path = args.data_path
    if args.output_dir is not None:
        cfg.output_dir = args.output_dir
    if args.device is not None:
        cfg.device = args.device
    if args.dtype is not None:
        cfg.dtype = args.dtype
    if args.epochs is not None:
        cfg.epochs = (args.epochs,)
        cfg.learning_rates = (args.lr if args.lr is not None else cfg.learning_rates[0],)
    elif args.lr is not None:
        cfg.learning_rates = tuple(args.lr for _ in cfg.learning_rates)
    if args.batch_size is not None:
        cfg.batch_size = args.batch_size
        cfg.num_batches = None
    if args.num_batches is not None:
        cfg.num_batches = args.num_batches
        cfg.batch_size = None
    if args.hidden_width is not None or args.hidden_depth is not None:
        width = args.hidden_width if args.hidden_width is not None else cfg.hidden_layers[0]
        depth = args.hidden_depth if args.hidden_depth is not None else len(cfg.hidden_layers)
        cfg.hidden_layers = (width,) * depth
    if args.seed is not None:
        cfg.seed = args.seed
    if args.parameterized is not None:
        cfg.parameterized = args.parameterized
    train(cfg, resume=args.resume)


if __name__ == "__main__":
    main()
\end{lstlisting}

\bigskip
\noindent{\large\bfseries evaluate.py}
\nopagebreak
\begin{lstlisting}
from __future__ import annotations

import argparse
from pathlib import Path

import numpy as np
import torch

from .config import PhysicsConfig
from .data import radius_from_h5_path
from .model import TwoPhasePINN
from .visualize import error_metrics, plot_field_comparison, predict_cfd_grid


def mesh_inputs(x: np.ndarray, y: np.ndarray, t: np.ndarray, radius: float | None = None) -> torch.Tensor:
    yy, tt, xx = np.meshgrid(y, t, x)
    columns = [xx.ravel(), yy.ravel(), tt.ravel()]
    if radius is not None:
        columns.append(np.full(xx.size, radius, dtype=np.float64))
    return torch.from_numpy(np.stack(columns, axis=1).astype(np.float32))


def reshape_prediction(values: torch.Tensor, x: np.ndarray, y: np.ndarray, t: np.ndarray) -> np.ndarray:
    return values.detach().cpu().numpy().reshape(len(t), len(y), len(x), order="C")


def load_model(checkpoint_path: Path, device: torch.device) -> tuple[TwoPhasePINN, dict]:
    checkpoint = torch.load(checkpoint_path, map_location=device)
    cfg = checkpoint["config"]
    physics = PhysicsConfig(**cfg["physics"])
    input_dim = int(checkpoint["model_state"]["net.trunk.0.weight"].shape[1])
    model = TwoPhasePINN(
        tuple(cfg["hidden_layers"]),
        physics,
        cfg.get("activation", "tanh"),
        tuple(cfg["loss_weights_pde"]),
        input_dim=input_dim,
    )
    dtype = torch.float64 if cfg.get("dtype") == "float64" else torch.float32
    model.to(device=device, dtype=dtype)
    model.load_state_dict(checkpoint["model_state"])
    model.eval()
    return model, cfg


def main() -> None:
    parser = argparse.ArgumentParser(description="Evaluate a trained rising bubble PINN checkpoint.")
    parser.add_argument("--checkpoint", type=Path, required=True)
    parser.add_argument("--data-path", type=Path, required=True)
    parser.add_argument("--device", default="auto")
    parser.add_argument("--start", type=int, default=0)
    parser.add_argument("--end", type=int, default=151)
    parser.add_argument("--temporal-step", type=int, default=10)
    parser.add_argument("--spatial-step", type=int, default=4)
    parser.add_argument("--time-index", type=int, default=-1)
    parser.add_argument("--save-figure", type=Path)
    args = parser.parse_args()

    pred, true, coords = predict_cfd_grid(
        args.checkpoint,
        args.data_path,
        device=args.device,
        start=args.start,
        end=args.end,
        temporal_step=args.temporal_step,
        spatial_step=args.spatial_step,
    )
    idx = args.time_index
    print(f"Predicted arrays at resolution t={len(coords['t'])}, y={len(coords['y'])}, x={len(coords['x'])}")
    print(f"Radius: {radius_from_h5_path(args.data_path):.4g}")
    print(f"u range: {pred['u'].min():.4e} .. {pred['u'].max():.4e}")
    print(f"v range: {pred['v'].min():.4e} .. {pred['v'].max():.4e}")
    print(f"p range: {pred['p'].min():.4e} .. {pred['p'].max():.4e}")
    print(f"alpha range: {pred['alpha'].min():.4e} .. {pred['alpha'].max():.4e}")
    print(error_metrics(pred, true).to_string(index=False))

    fig = plot_field_comparison(pred, true, coords, time_index=idx, save_path=args.save_figure)
    if not args.save_figure:
        import matplotlib.pyplot as plt

        plt.show()


if __name__ == "__main__":
    main()
\end{lstlisting}

\bigskip
\noindent{\large\bfseries visualize.py}
\nopagebreak
\begin{lstlisting}
from __future__ import annotations

from pathlib import Path

import matplotlib.pyplot as plt
import numpy as np
import pandas as pd
import torch

from .config import PhysicsConfig
from .data import load_cfd, radius_from_h5_path
from .model import TwoPhasePINN


def resolve_device(device: str = "auto") -> torch.device:
    if device == "auto":
        return torch.device("cuda" if torch.cuda.is_available() else "cpu")
    requested = torch.device(device)
    if requested.type == "cuda" and not torch.cuda.is_available():
        raise RuntimeError("CUDA was requested, but torch.cuda.is_available() is False.")
    return requested


def _mesh_inputs(x: np.ndarray, y: np.ndarray, t: np.ndarray, radius: float | None = None) -> torch.Tensor:
    tt, yy, xx = np.meshgrid(t, y, x, indexing="ij")
    columns = [xx.ravel(), yy.ravel(), tt.ravel()]
    if radius is not None:
        columns.append(np.full(xx.size, radius, dtype=np.float64))
    return torch.from_numpy(np.stack(columns, axis=1).astype(np.float32))


def _reshape(values: torch.Tensor, t: np.ndarray, y: np.ndarray, x: np.ndarray) -> np.ndarray:
    return values.detach().cpu().numpy().reshape(len(t), len(y), len(x))


def _to_yx(values: np.ndarray) -> np.ndarray:
    return np.transpose(values, (0, 2, 1))


def load_model(checkpoint_path: Path, device: torch.device) -> tuple[TwoPhasePINN, dict]:
    checkpoint = torch.load(checkpoint_path, map_location=device)
    cfg = checkpoint["config"]
    physics = PhysicsConfig(**cfg["physics"])
    input_dim = int(checkpoint["model_state"]["net.trunk.0.weight"].shape[1])
    model = TwoPhasePINN(
        tuple(cfg["hidden_layers"]),
        physics,
        cfg.get("activation", "tanh"),
        tuple(cfg["loss_weights_pde"]),
        input_dim=input_dim,
    )
    dtype = torch.float64 if cfg.get("dtype") == "float64" else torch.float32
    model.to(device=device, dtype=dtype)
    model.load_state_dict(checkpoint["model_state"])
    model.eval()
    return model, cfg


def predict_cfd_grid(
    checkpoint_path: Path,
    data_path: Path,
    device: str = "auto",
    start: int = 0,
    end: int = 151,
    temporal_step: int = 10,
    spatial_step: int = 4,
    batch_size: int = 262_144,
) -> tuple[dict[str, np.ndarray], dict[str, np.ndarray], dict]:
    device_obj = resolve_device(device)
    model, cfg = load_model(Path(checkpoint_path), device_obj)
    physics = PhysicsConfig(**cfg["physics"])
    dtype = torch.float64 if cfg.get("dtype") == "float64" else torch.float32

    cfd = load_cfd(Path(data_path), start=start, end=end, temporal_step=temporal_step, spatial_step=spatial_step)
    x = cfd["x"] / physics.l_ref
    y = cfd["y"] / physics.l_ref
    t = cfd["time"] / physics.l_ref
    radius = radius_from_h5_path(Path(data_path)) / physics.l_ref
    input_dim = int(model.net.trunk[0].weight.shape[1])
    xyt = _mesh_inputs(x, y, t, radius if input_dim == 4 else None).to(device=device_obj, dtype=dtype)

    with torch.no_grad():
        u, v, p, alpha = model.predict(xyt, batch_size=batch_size)

    pred = {
        "u": _reshape(u, t, y, x),
        "v": _reshape(v, t, y, x),
        "p": _reshape(p, t, y, x),
        "alpha": _reshape(alpha, t, y, x),
    }
    true = {
        "u": _to_yx(cfd["velocity_x"]),
        "v": _to_yx(cfd["velocity_y"]),
        "p": _to_yx(cfd["pressure"]) / (physics.p_ref_rho * physics.u_ref**2),
        "alpha": (_to_yx(cfd["levelset"]) < 0.0).astype(np.float32),
    }
    coords = {
        "x": x,
        "y": y,
        "t": t,
        "time_raw": cfd["time"],
        "radius": radius,
        "physics": physics,
        "config": cfg,
    }
    return pred, true, coords


def error_metrics(pred: dict[str, np.ndarray], true: dict[str, np.ndarray]) -> pd.DataFrame:
    rows = []
    for name in ("u", "v", "p", "alpha"):
        diff = pred[name] - true[name]
        rmse = float(np.sqrt(np.mean(diff**2)))
        mae = float(np.mean(np.abs(diff)))
        denom = float(np.sqrt(np.mean(true[name] ** 2)))
        rows.append(
            {
                "field": name,
                "rmse": rmse,
                "mae": mae,
                "relative_rmse": rmse / denom if denom > 0 else np.nan,
                "pred_min": float(np.min(pred[name])),
                "pred_max": float(np.max(pred[name])),
                "true_min": float(np.min(true[name])),
                "true_max": float(np.max(true[name])),
            }
        )
    return pd.DataFrame(rows)


def align_pressure_gauge(pred: dict[str, np.ndarray], true: dict[str, np.ndarray]) -> dict[str, np.ndarray]:
    aligned = dict(pred)
    aligned["p"] = pred["p"] + (true["p"][:, -1:, :].mean() - pred["p"][:, -1:, :].mean())
    return aligned


def plot_field_comparison(
    pred: dict[str, np.ndarray],
    true: dict[str, np.ndarray],
    coords: dict,
    time_index: int = -1,
    fields: tuple[str, ...] = ("alpha", "u", "v", "p"),
    cmap: str = "viridis",
    error_cmap: str = "coolwarm",
    align_pressure: bool = True,
    save_path: Path | None = None,
):
    x = coords["x"]
    y = coords["y"]
    t_value = coords["time_raw"][time_index]
    fig, axes = plt.subplots(len(fields), 3, figsize=(13, 3.2 * len(fields)), constrained_layout=True)
    if len(fields) == 1:
        axes = axes.reshape(1, 3)

    for row, field in enumerate(fields):
        pred_field = pred[field][time_index]
        true_field = true[field][time_index]
        if field == "p" and align_pressure:
            pred_field = pred_field + (true_field[-1, :].mean() - pred_field[-1, :].mean())
        err_field = pred_field - true_field
        vmin = min(float(pred_field.min()), float(true_field.min()))
        vmax = max(float(pred_field.max()), float(true_field.max()))
        err_abs = float(np.max(np.abs(err_field)))

        images = [
            axes[row, 0].contourf(x, y, true_field, levels=60, cmap=cmap, vmin=vmin, vmax=vmax),
            axes[row, 1].contourf(x, y, pred_field, levels=60, cmap=cmap, vmin=vmin, vmax=vmax),
            axes[row, 2].contourf(x, y, err_field, levels=60, cmap=error_cmap, vmin=-err_abs, vmax=err_abs),
        ]
        axes[row, 0].set_ylabel(field)
        axes[row, 0].set_title(f"CFD, t={t_value:.4g}")
        axes[row, 1].set_title("PINN" if field != "p" or not align_pressure else "PINN, aligned")
        axes[row, 2].set_title("PINN - CFD")

        for ax, image in zip(axes[row], images):
            ax.set_aspect("equal")
            ax.set_xlabel("x / L")
            fig.colorbar(image, ax=ax)

    if save_path is not None:
        save_path = Path(save_path)
        save_path.parent.mkdir(parents=True, exist_ok=True)
        fig.savefig(save_path, dpi=180)
    return fig


def plot_centerline(
    pred: dict[str, np.ndarray],
    true: dict[str, np.ndarray],
    coords: dict,
    field: str = "alpha",
    time_index: int = -1,
    x_value: float = 0.0,
):
    x = coords["x"]
    y = coords["y"]
    x_idx = int(np.argmin(np.abs(x - x_value)))
    fig, ax = plt.subplots(figsize=(6, 4), constrained_layout=True)
    ax.plot(true[field][time_index, :, x_idx], y, label="CFD")
    ax.plot(pred[field][time_index, :, x_idx], y, label="PINN")
    ax.set_xlabel(field)
    ax.set_ylabel("y / L")
    ax.set_title(f"{field} centerline at x={x[x_idx]:.3g}")
    ax.grid(True)
    ax.legend()
    return fig
\end{lstlisting}
